% Inline source for IEEEtran. Requires graphicx, amsmath and american circuitikz.
\section{Transistor Implementation of the AFC}
\label{sec:r1afc}
The R1 AFC compares edges of the 100-MHz reference and the physical
divide-by-32 feedback. A higher five-bit code increases the ring frequency.
The controller starts at code 16 and can test four further candidates
using steps of 8, 4, 2, and 1. Each candidate reaches the detector before
the controller can hand it to the fine loop.

\begin{figure*}[t]
\centering\resizebox{0.98\textwidth}{!}{%
\begin{circuitikz}[american,x=1cm,y=1cm,font=\small,
 blk/.style={draw,rounded corners=1.5pt,align=center,minimum width=2.5cm,
 minimum height=.76cm,inner sep=4pt}, ar/.style={->,>=latex}]
  \node[blk] (ep) at (1.7,2.2) {First FB edge\\sample REF level};
  \node[blk] (sel) at (5.1,2.2) {Phase selector\\REF or inverted REF};
  \node[blk] (count) at (8.8,2.2) {Lead / lag counters\\8-bit Gray, init.\ 1 / 0};
  \node[blk] (fbg) at (1.7,.55) {FB gate\\enable latched low};
  \node[blk] (div) at (5.1,.55) {FB counter\\8-bit Gray, init.\ 0};
  \node[blk] (sync) at (12.3,1.38) {Two DFFs per bit\\then Gray decode};
  \node[blk] (cmp) at (12.3,-.7) {Compare counts\\qualify direction};
  \node[blk] (snap) at (8.8,-.7) {Decision snapshot\\$3\times8$ bits + direction};
  \node[blk] (fsm) at (5.1,-.7) {Trial controller\\state and timers};
  \node[blk] (code) at (1.7,-.95) {Five-bit band register\\steps $8,4,2,1$};
  \draw[ar] (-.95,2.2) node[left]{REF} -- (ep.west);
  \draw[ar] (-.95,.55) node[left]{FB} -- (fbg.west);
  \draw[ar] (ep.east) -- (sel.west); \draw[ar] (sel.east) -- (count.west);
  \draw[ar] (fbg.east) -- (div.west);
  \draw[ar] (3.26,.55) node[circ]{} -- (3.26,1.32) -| (ep.south);
  \draw[ar] (count.east) -- ++(.20,0) |- ([yshift=2.5mm]sync.west);
  \draw[ar] (div.east) -- (10.35,.55) |- ([yshift=-2.5mm]sync.west);
  \draw[ar] (sync.south) -- (cmp.north); \draw[ar] (cmp.west) -- (snap.east);
  \draw[ar] (snap.west) -- (fsm.east); \draw[ar] (fsm.west) -- ++(-.27,0) |- (code.east);
  \draw[ar] (code.west) -- ++(-.85,0) node[left]{$S[4{:}0]$};
  \draw[ar] (fsm.south) -- ++(0,-.54) node[below]{DONE, FAIL, HOLD, FineEnable, PFD reset};
\end{circuitikz}}
\caption{R1 frequency detector and search controller. The phase selector also receives REF and its complement. The gated FB clocks both the FB counter and the first-edge sample. The synchronized counts share the controller clock. Trial control supplies the counter enables and resets.}
\label{fig:r1afc-datapath}
\end{figure*}

The first admitted FB edge samples the reference level and assigns
the two reference phases to lead and lag counters (Fig.~\ref{fig:r1afc-datapath}).
The lead counter starts with one credit. A chain of 16 inverters guards
reference-counter reset release after that first edge. The controller
samples each Gray word through two flip-flops per bit, then decodes it.
With lead, lag, and divided-feedback counts $C_L,C_G,C_D$, respectively,
it qualifies a slow decision when $C_G-C_D\geq 2$ and a fast decision
when $C_D-C_L\geq 2$. The unsigned subtractors check the carry first.
Two consecutive observations in the same direction qualify a request.
Reaching lag count 128 without a qualified direction ends a trial as near.

\begin{figure}[t]
\centering\resizebox{.93\columnwidth}{!}{%
\begin{circuitikz}[american,x=1cm,y=1cm,font=\small,
 s/.style={draw,rounded corners=1.5pt,minimum width=1.72cm,
 minimum height=.60cm,inner sep=3pt,align=center}, a/.style={->,>=latex}]
 \node[s] (set) at (1,1.5) {SETTLE}; \node[s] (arm) at (4.8,1.5) {ARM};
 \node[s] (meas) at (4.8,0) {MEASURE}; \node[s] (dr) at (1,0) {DRAIN};
 \node[s] (ok) at (1,-1.8) {DONE}; \node[s] (fail) at (4.8,-1.8) {FAIL};
 \draw[a] (1,2.32) node[above]{START} -- (set.north);
 \draw[a] (set.east) -- node[above]{4 REF} (arm.west);
 \draw[a] (arm.south) -- node[right,font=\scriptsize]{epoch valid} (meas.north);
 \draw[a] (meas.west) -- node[above]{decision} (dr.east);
 \draw[a] (dr.west) -- ++(-.45,0) |- ([yshift=-1mm]set.west);
 \node[rotate=90,anchor=south,font=\scriptsize] at (-.33,.75) {next code};
 \draw[a] (dr.south) -- node[left,font=\scriptsize]{accept} (ok.north);
 \draw[a] (meas.south) -- node[right,font=\scriptsize]{watchdog / conflict} (fail.north);
 \draw[a] ([yshift=1mm]arm.east) -- ++(1.62,0) |- (fail.east);
 \node[rotate=90,anchor=south,font=\scriptsize] at (7.15,.5) {ARM timeout};
 \draw[a] ([xshift=4mm]dr.south) -- ++(0,-.91) -| ([xshift=-5mm]fail.north);
 \node[anchor=south,font=\scriptsize] at (2.95,-1.2) {endpoint failure};
\end{circuitikz}}
\caption{R1 state transitions. DRAIN holds the captured counts for two controller cycles. A terminal decision accepts a near result or the last tested candidate. An outward direction at code 1 or 31 on a terminal decision sets FAIL.}
\label{fig:r1afc-fsm}
\end{figure}

An empty ARM window times out after 32 reference clocks. MEASURE has
a 192-clock watchdog. The circuit stores the three decision words and
direction before ending the measurement. On success, two further
reference clocks separate DONE from FineEnable. HOLD is the complement
of FineEnable. The PFD reset is asserted when reset is high or FineEnable
is low.

\subsection{Static CMOS Cells}
The generator expands the gates into one flat Spectre subcircuit.
Figure~\ref{fig:r1afc-cells} shows the transistor connections of the
gate constructors. All devices use $L=0.13~\mu\mathrm{m}$. NMOS bulk
connects to VSS and PMOS bulk connects to VDD. Repeated
input labels in a panel denote the same electrical net.

\begin{figure*}[t]
\centering\resizebox{.97\textwidth}{!}{%
\begin{circuitikz}[american,x=1cm,y=1cm,font=\small]
\begin{scope}
 \node[pmos] (pi) at (0,2.1) {}; \node[nmos] (ni) at (0,.55) {};
 \draw (pi.S) -- (0,2.9) node[above]{VDD};
 \draw (ni.S) -- (0,-.2) node[below]{VSS};
 \draw (pi.D) -- (ni.D); \draw (0,1.325) node[circ]{} -- ++(.65,0) node[right]{$Y$};
 \draw (pi.G) -- (-.84,2.1) -- (-.84,.55) -- (ni.G);
 \draw (-.84,1.325) node[circ]{} -- ++(-.35,0) node[left]{$A$};
 \node at (0,-1.64) {(a) INV};
\end{scope}
\begin{scope}[shift={(3.55,0)}]
 \node[pmos] (pa) at (-.55,2.1) {}; \node[pmos,xscale=-1] (pb) at (.55,2.1) {};
 \node[nmos] (na) at (0,.75) {}; \node[nmos] (nb) at (0,-.28) {};
 \draw (pa.S) -- (-.55,2.92) -- (.55,2.92) -- (pb.S);
 \node[above] at (0,2.92) {VDD};
 \draw (pa.D) -- (-.55,1.42) -- (.55,1.42) -- (pb.D);
 \draw (0,1.42) node[circ]{} -- (na.D); \draw (na.S) -- (nb.D);
 \draw (.55,1.42) -- ++(.55,0) node[right]{$Y$};
 \draw (nb.S) -- (0,-.99) node[below]{VSS};
 \draw (pa.G) -- ++(-.12,0) node[left]{$A$}; \draw (pb.G) -- ++(.12,0) node[right]{$B$};
 \draw (na.G) -- ++(-.12,0) node[left]{$A$}; \draw (nb.G) -- ++(-.12,0) node[left]{$B$};
 \node at (0,-1.64) {(b) NAND2};
\end{scope}
\begin{scope}[shift={(7.6,0)}]
 \node[pmos] (pb) at (0,2.35) {}; \node[pmos] (pa) at (0,1.32) {};
 \node[nmos] (na) at (-.55,-.05) {}; \node[nmos,xscale=-1] (nb) at (.55,-.05) {};
 \draw (pb.S) -- (0,3.02) node[above]{VDD}; \draw (pb.D) -- (pa.S);
 \draw (pa.D) -- (0,.64) node[circ]{}; \draw (na.D) -- (-.55,.64) -- (.55,.64) -- (nb.D);
 \draw (.55,.64) -- ++(.55,0) node[right]{$Y$};
 \draw (na.S) -- (-.55,-.81) -- (.55,-.81) -- (nb.S); \node[below] at (0,-.81) {VSS};
 \draw (pa.G) -- ++(-.12,0) node[left]{$A$}; \draw (pb.G) -- ++(-.12,0) node[left]{$B$};
 \draw (na.G) -- ++(-.12,0) node[left]{$A$}; \draw (nb.G) -- ++(.12,0) node[right]{$B$};
 \node at (0,-1.64) {(c) NOR2};
\end{scope}
\begin{scope}[shift={(11.6,0)}]
 \foreach \i/\gate in {0/A,1/B,2/C,3/D}{
   \node[pmos] (p\i) at ({1.02*\i},1.45) {};
   \draw (p\i.S) -- ({1.02*\i},2.02); \draw (p\i.D) -- ({1.02*\i},.85);
   \draw (p\i.G) -- ++(-.05,0) node[left,font=\scriptsize]{$\gate$};
 }
 \draw (0,2.02) -- (3.06,2.02); \node[above] at (1.53,2.02) {VDD};
 \draw (0,.85) -- (3.06,.85); \node[below] at (1.53,.85) {$Y$};
 \foreach \i/\gate in {0/A,1/B,2/C,3/D}{
   \node[nmos] (m\i) at (4.30,{2.1-.9*\i}) {};
   \draw (m\i.G) -- ++(-.05,0) node[left,font=\scriptsize]{$\gate$};
 }
 \draw (m0.D) -- ++(0,.12) node[above]{$Y$};
 \draw (m0.S) -- (m1.D); \draw (m1.S) -- (m2.D); \draw (m2.S) -- (m3.D);
 \draw (m3.S) -- ++(0,-.12) node[below]{VSS};
 \node at (1.95,-1.64) {(d) NAND4};
\end{scope}
\begin{scope}[shift={(19.0,0)}]
 \node[pmos,rotate=-90] (pg) at (0,1.8) {};
 \node[nmos,rotate=90] (ng) at (0,.5) {};
 \draw (pg.D) -- (-.85,1.8) -- (-.85,.5) -- (ng.D);
 \draw (pg.S) -- (.85,1.8) -- (.85,.5) -- (ng.S);
 \draw (-.85,1.13) node[circ]{} -- ++(-.3,0) node[left]{$X$};
 \draw (.85,1.13) node[circ]{} -- ++(.3,0) node[right]{$Z$};
 \draw (ng.G) -- ++(0,-.14) node[below]{$E$};
 \draw (pg.G) -- ++(0,.14) node[above]{$\overline E$};
 \node at (0,-1.64) {(e) TG};
\end{scope}
\end{circuitikz}}
\caption{Distinct gate topologies of the R1 AFC. For the NAND4 the two networks share output $Y$. The transmission gate (TG) connects $X$ and $Z$ when $E=1$. The latch adds the reset NMOS shown in Fig.~\ref{fig:r1afc-latch}.}
\label{fig:r1afc-cells}
\end{figure*}

The default inverter has $W_n/W_p=0.8/2.4~\mu\mathrm{m}$. Its latch and
clock-driver variants use $0.5/1.5$ and $8/24~\mu\mathrm{m}$.
NAND2, NOR2 and NAND4 use per-device widths of $1.6/2.4$,
$0.8/4.8$ and $3.2/2.4~\mu\mathrm{m}$, respectively. The TG uses
$1.2/3.6~\mu\mathrm{m}$. For an XOR, the generator forms
$U=\overline{AB}$ and then
$Y=\overline{\overline{AU}\,\overline{BU}}$. A two-input multiplexer
uses NAND2 cells and a complemented select. These same cells form the
Gray converters, incrementers, difference detectors and code update logic.

\begin{figure}[t]
\centering\resizebox{.98\columnwidth}{!}{%
\begin{circuitikz}[american,x=1cm,y=1cm,font=\PLLfigurefont\small,
 c/.style={draw,minimum width=1.12cm,minimum height=.66cm,align=center},
 ar/.style={->,>=latex}]
 \node[c] (di) at (1.05,1.05) {TG}; \node[c] (ii) at (3.65,1.05) {INV};
 \node[c] (io) at (5.5,1.05) {INV}; \node[c] (fb) at (3.65,2.7) {TG};
 \draw (-.37,1.05) node[left]{$D$} -- (di.west);
 \draw (di.east) -- (ii.west); \draw (ii.east) -- (io.west);
 \draw (io.east) -- (6.67,1.05) node[right]{$Q$};
 \draw (6.32,1.05) node[circ]{} |- (fb.east); \draw (fb.west) -| (2.32,1.05) node[circ]{};
 \draw (di.north) -- ++(0,.28) node[above]{open};
 \draw (fb.north) -- ++(0,.16) node[above]{hold};
 \draw (4.58,1.05) node[circ]{} -- ++(0,.32) node[above]{$\overline Q$};
 \node[nmos] (mr) at (2.32,-.12) {};
 \draw (mr.D) -- (2.32,1.05); \draw (mr.G) -- ++(-.14,0) node[left]{RST};
 \draw (mr.S) -- ++(0,-.12) node[below]{VSS};
 \node[anchor=west,font=\PLLfigurefont\scriptsize] at (2.45,.6) {$X$};
 \node[anchor=west,font=\PLLfigurefont\scriptsize] at (2.65,-.1) {$W_n=2~\mu$m};
 \begin{scope}[shift={(0,-2.55)}]
 \node[c] (ml) at (1.5,0) {Master latch};
 \node[c] (sl) at (4.8,0) {Slave latch};
 \draw[ar] (-.35,0) node[left]{$D$} -- (ml.west);
 \draw[ar] (ml.east) -- node[above]{$Q_M$} (sl.west);
 \draw[ar] (sl.east) -- (6.65,0) node[right]{$Q$};
 \node[below,align=center] at (ml.south) {open $L$\\hold $H$};
 \node[below,align=center] at (sl.south) {open $H$\\hold $L$};
 \end{scope}
\end{circuitikz}}
\caption{R1 resettable static latch (top) and ordinary positive-edge DFF (bottom). Each TG receives an enable and its complement. The feedback TG returns $Q$ to storage node $X$ during hold. Active reset pulls $X$ to VSS in each latch.}
\label{fig:r1afc-latch}
\end{figure}

Two latches form one DFF. The preset variant inverts $D$ and exposes
the slave's complementary output. Each shared clock group forms $K_d$
with six cascaded inverters. With active-high reset $R$, the interlock
computes
\begin{align}
 L &= \overline K\,\overline{K_d}\,\overline H\,\overline R, &
 H &= K\,K_d\,\overline L\,\overline R .
\end{align}
A NAND4 and two large inverters drive each phase and its complement.

R1 contains 146 DFFs and one additional latch on the feedback clock
gate. The 293 latches use 586 TGs, 586 latch inverters and 293 reset
NMOS devices. Outside these latches, the generator emits 531 inverters
(including clock drivers), 740 NAND2s, 79 NOR2s and 20 NAND4s. The
total is $293(9)+531(2)+740(4)+79(4)+20(8)=7{,}135$ MOS instances
in the AFC core. These counts can be checked against the selected R1
gate generator, its port manifest and the generated transistor netlist.

